\section*{Chapitre \ref{ch:g11:cancer} --- Atteintes du génome et cancer}

\begin{solution}{exo:g11:cancer:1}
\omterm{def:g11:cancer:cancer}{Cancer}~: des \omterm{def:g10:cells-common-unit:cell}{cellules} qui se divisent sans les contrôles normaux et
envahissent le tissu voisin. \omterm{def:g11:cancer:cancer}{Tumeur}~: la masse formée par un tel clone.
Métastase~: une nouvelle \omterm{def:g11:cancer:cancer}{tumeur} fondée ailleurs par des \omterm{def:g10:cells-common-unit:cell}{cellules} qui
ont voyagé par le sang ou la lymphe.
\end{solution}

\begin{solution}{exo:g11:cancer:2}
Un \omterm{def:g10:universal-dna:gene}{gène} dont la \omterm{def:g11:gene-expression:protein}{protéine} relaie les signaux de croissance et pousse la
\omterm{def:g10:cells-common-unit:cell}{cellule} dans le cycle~; muté de sorte que la \omterm{def:g11:gene-expression:protein}{protéine} soit constamment
active, il devient un oncogène. Un accélérateur bloqué entraîne la
\omterm{def:g10:cells-common-unit:cell}{cellule} quoi que fasse l'autre \omterm{def:g10:universal-dna:gene}{allèle}.
\end{solution}

\begin{solution}{exo:g11:cancer:3}
Elle arrête le cycle quand l'\omterm{def:g10:universal-dna:information}{ADN} est abîmé et déclenche la réparation
ou l'autodestruction. Un seul \omterm{def:g10:universal-dna:gene}{allèle} fonctionnel fabrique encore assez
de \omterm{def:g11:gene-expression:protein}{protéine} pour freiner~; seule la perte des deux ôte le frein.
\end{solution}

\begin{solution}{exo:g11:cancer:4}
Le rayonnement ultraviolet (\omterm{def:g11:cancer:cancer}{cancers} de la peau), les goudrons du tabac
(\omterm{def:g11:cancer:cancer}{cancer} du poumon), les rayonnements ionisants, l'amiante, l'alcool, le
papillomavirus (\omterm{def:g11:cancer:cancer}{cancer} du col de l'utérus).
\end{solution}

\begin{solution}{exo:g11:cancer:5}
Toutes ses \omterm{def:g10:cells-common-unit:cell}{cellules} descendent d'une \omterm{def:g10:cells-common-unit:cell}{cellule} où la première \omterm{def:g11:mutations:mutation}{mutation}
est survenue~; elles partagent cette \omterm{def:g11:mutations:mutation}{mutation} et celles qui ont suivi,
ainsi que des marques comme le même \omterm{def:g11:cell-cycle-mitosis:chromosome}{chromosome} X inactif.
\end{solution}

\begin{solution}{exo:g11:cancer:6}
Environ 200 à 40 ans et 2700 à 80 ans~: un facteur 13 pour un facteur 2
sur l'âge --- plusieurs événements qui se multiplient, et non une cause
unique proportionnelle.
\end{solution}

\begin{solution}{exo:g11:cancer:7}
Environ 10 et 17. Un fumeur de 20 par jour~: $17 \times 1\% = 17\%$,
soit un sur six.
\end{solution}

\begin{solution}{exo:g11:cancer:8}
Sans p53, la \omterm{def:g10:cells-common-unit:cell}{cellule} ne marque plus de pause pour réparer son \omterm{def:g10:universal-dna:information}{ADN} abîmé
avant la réplication, et ne se détruit plus quand les dégâts sont
graves~: les lésions sont copiées en \omterm{def:g11:mutations:mutation}{mutations} à chaque division, et
des \omterm{def:g10:cells-common-unit:cell}{cellules} au génome bouleversé survivent.
\end{solution}

\begin{solution}{exo:g11:cancer:9}
Chaque \omterm{def:g10:cells-common-unit:cell}{cellule} de son corps est déjà privée d'un \omterm{def:g10:universal-dna:gene}{allèle}~; une seule
\omterm{def:g11:mutations:mutation}{mutation} somatique de l'autre achève la perte, au lieu de deux
\omterm{def:g11:mutations:mutation}{mutations} indépendantes. Une étape étant franchie dès la naissance, les
étapes restantes sont atteintes plus tôt et dans un plus grand nombre
de \omterm{def:g10:cells-common-unit:cell}{cellules}. Ses sœurs héritent l'\omterm{def:g10:universal-dna:gene}{allèle} du même parent avec une
probabilité de $1/2$.
\end{solution}

\begin{solution}{exo:g11:cancer:10}
Les médicaments tuent les \omterm{def:g10:cells-common-unit:cell}{cellules} en division sans discernement~: les
\omterm{def:g10:cells-common-unit:cell}{cellules} des racines des cheveux, la paroi intestinale et la moelle
osseuse se divisent vite et sont frappées en même temps que la \omterm{def:g11:cancer:cancer}{tumeur}
--- d'où la chute des cheveux, les nausées et le manque de globules
blancs.
\end{solution}

\begin{solution}{exo:g11:cancer:11}
Arrêter met fin à l'apport de \omterm{def:g11:mutations:mutation}{mutations} nouvelles, et les \omterm{def:g10:cells-common-unit:cell}{cellules} de
la paroi sont peu à peu remplacées, ce qui élimine bien des clones
\omterm{def:g11:genes-and-disease:genetic}{porteurs} des premières étapes. Mais certains clones à plusieurs étapes
persistent, et les \omterm{def:g11:mutations:mutation}{mutations} déjà présentes ne s'effacent pas~; le
risque reste au-dessus de celui d'un non-fumeur.
\end{solution}

\begin{solution}{exo:g11:cancer:12}
$10^{12} \times 10^{-7} = 10^5$ \omterm{def:g10:cells-common-unit:cell}{cellules} acquièrent la première
\omterm{def:g11:mutations:mutation}{mutation}. Chacune fonde un clone qui peut se diviser des milliers de
fois, si bien que la deuxième \omterm{def:g11:mutations:mutation}{mutation} n'est pas cherchée dans une
\omterm{def:g10:cells-common-unit:cell}{cellule} mais dans toutes les divisions du clone, et de même pour la
troisième~: les étapes multiplient la cible, et la probabilité réelle
est de très nombreux ordres de grandeur au-dessus de $10^{-21}$.
\end{solution}

\begin{solution}{exo:g11:cancer:13}
La \omterm{def:g11:gene-expression:protein}{protéine} virale ôte le frein sans toucher au \omterm{def:g10:universal-dna:gene}{gène}~: la \omterm{def:g10:cells-common-unit:cell}{cellule}
infectée se comporte comme si p53 était perdue et accumule les
\omterm{def:g11:mutations:mutation}{mutations} qui achèvent le processus. Le vaccin fait détruire le virus
par le système immunitaire avant qu'il n'infecte le col~; une fois les
\omterm{def:g10:cells-common-unit:cell}{cellules} infectées et transformées, le vaccin ne sert plus --- d'où la
vaccination au début de l'adolescence.
\end{solution}

\begin{solution}{exo:g11:cancer:14}
Un polype est un clone qui a franchi une ou deux étapes~; le retirer
supprime les \omterm{def:g10:cells-common-unit:cell}{cellules} dans lesquelles les étapes restantes pourraient
survenir, si bien que le \omterm{def:g11:cancer:cancer}{cancer} ne se forme jamais. Comme la suite
polype--cancer prend dix à vingt ans, un intervalle de dix ans attrape
la plupart des polypes avant qu'ils ne l'achèvent.
\end{solution}

\begin{solution}{exo:g11:cancer:15}
L'incidence monte avec l'âge parce que les étapes prennent des
décennies, mais ces étapes sont franchies au cours de la vie, en grande
partie sous l'effet d'expositions --- le tabac à lui seul rend compte
d'un tiers des décès par \omterm{def:g11:cancer:cancer}{cancer}, et les infections, l'alcool, le soleil
et l'alimentation d'une bonne part du reste. Supprimer une exposition à
tout âge réduit les \omterm{def:g11:mutations:mutation}{mutations} encore à venir~; le dépistage supprime
des clones avant la dernière étape. La vieillesse est le moment où la
note est présentée, non celui où elle est écrite.
\end{solution}

\begin{solution}{pb:g11:cancer:1}
\textbf{1.} Elles sont présentes dans la \omterm{def:g11:cancer:cancer}{tumeur} et absentes des autres
\omterm{def:g10:cells-common-unit:cell}{cellules} du patient~: elles sont donc apparues dans une \omterm{def:g10:cells-common-unit:cell}{cellule} du
corps. Des \omterm{def:g11:mutations:mutation}{mutations} germinales seraient dans chaque \omterm{def:g10:cells-common-unit:cell}{cellule}, tissu
sain compris.

\textbf{2.} Que la plupart des \omterm{def:g11:mutations:mutation}{mutations} de la \omterm{def:g11:cancer:cancer}{tumeur} ont été causées
par le benzopyrène --- l'\omterm{def:g11:mutations:mutagen}{agent mutagène} de la fumée de tabac y a laissé
sa signature.

\textbf{3.} Une \omterm{def:g11:mutations:mutation}{mutation} non-sens~: un \omterm{def:g11:gene-expression:code}{codon} d'arrêt prématuré tronque
p53, qui ne peut plus fonctionner.

\textbf{4.} p53 est un frein~: un \omterm{def:g10:universal-dna:gene}{allèle} fonctionnel freine encore, et
il a donc fallu perdre les deux (l'un muté, l'autre délété). Le
récepteur est un accélérateur~: une seule \omterm{def:g11:gene-expression:protein}{protéine} constamment active
entraîne la \omterm{def:g10:cells-common-unit:cell}{cellule} quoi que fasse l'autre \omterm{def:g10:universal-dna:gene}{allèle}.

\textbf{5.} Le \omterm{def:g10:universal-dna:gene}{gène} du récepteur est l'accélérateur (\omterm{prop:g11:cancer:genes}{proto-oncogène}
devenu oncogène)~; p53 est le frein (suppresseur de \omterm{def:g11:cancer:cancer}{tumeur}).

\textbf{6.} $20 \times 365 \times 35 \approx \num{255000}$ cigarettes.

\textbf{7.} Environ $\num{255000}$ substitutions par \omterm{def:g10:cells-common-unit:cell}{cellule} ---
cinquante fois les 5000 trouvées. La plupart des lésions sont réparées,
et les \omterm{def:g10:cells-common-unit:cell}{cellules} trop mutées meurent~; 5000, c'est la fraction qui a
survécu.

\textbf{8.} Elles sont survenues tôt, dans la \omterm{def:g10:cells-common-unit:cell}{cellule} fondatrice ou ses
premières descendantes, si bien que chaque \omterm{def:g10:cells-common-unit:cell}{cellule} tumorale en a
hérité~; les autres sont apparues plus tard, dans des sous-clones.

\textbf{9.} $10^{10} \approx 2^{33}$~: 33 doublements, à 100 jours
chacun, soit environ 9 ans.

\textbf{10.} Si la \omterm{def:g11:cancer:cancer}{tumeur} a mis 9 ans ou plus à croître à partir de sa
\omterm{def:g10:cells-common-unit:cell}{cellule} fondatrice, la dernière \omterm{def:g11:mutations:mutation}{mutation} clé est survenue vers la
25\textsuperscript{e} année des 35 ans de tabagisme, ou plus tôt.

\textbf{11.} $17 \times 1\% = 17\%$.

\textbf{12.} 17 \omterm{def:g11:cancer:cancer}{cancers}, dont 1 serait survenu de toute façon~: 16 sur
100 sont attribuables au tabac.

\textbf{13.} Aucune \omterm{def:g11:mutations:mutation}{mutation} nouvelle ne s'ajoute~; la paroi se
renouvelle et bien des clones \omterm{def:g11:genes-and-disease:genetic}{porteurs} des premières étapes sont
éliminés et remplacés au fil des ans~; les clones restants ont moins
d'étapes qu'ils n'en auraient acquises en continuant.

\textbf{14.} Plusieurs \omterm{def:g11:mutations:mutation}{mutations} rares doivent s'accumuler dans une
même lignée, et la lignée doit ensuite atteindre une taille
détectable~: chaque étape prend des années, et leur somme fait des
décennies.

\textbf{15.} Une \omterm{def:g11:mutations:mutation}{mutation} par \omterm{def:g10:cells-common-unit:cell}{cellule} et par cigarette, dans quelque
$10^9$ \omterm{def:g10:cells-common-unit:cell}{cellules} de la paroi, sur $\num{255000}$ cigarettes, cela fait
$10^{14}$ \omterm{def:g11:mutations:mutation}{mutations} réparties dans le poumon~; quelques-unes toucheront
p53 ou un \omterm{prop:g11:cancer:genes}{proto-oncogène} dans une \omterm{def:g10:cells-common-unit:cell}{cellule} qui porte déjà une autre
étape. C'est le nombre d'essais qui rend l'improbable certain.

\textbf{16.} Fumeurs~: $\num{27000}$ décès pour $\num{11.25e6}$
personnes, soit environ 240 pour 100\,000~; non-fumeurs~: 3000 pour
$\num{33.75e6}$, soit environ 9 pour 100\,000.

\textbf{17.} Environ 27 fois --- du même ordre que les 17 à 25 de la
figure pour les gros fumeurs~; le chiffre de population mêle petits et
gros fumeurs et anciens fumeurs.

\textbf{18.} Environ 3000 par an, la part des non-fumeurs. Pas
immédiatement~: les \omterm{def:g11:mutations:mutation}{mutations} déjà présentes chez les fumeurs actuels
et anciens produiront des \omterm{def:g11:cancer:cancer}{cancers} encore vingt ou trente ans.

\textbf{19.} Dépister tous les fumeurs épargne environ 20~\% de
$\num{27000}$, soit quelque 5000 décès par an~; arrêter de fumer en
préviendrait à terme environ $\num{24000}$ --- cinq fois plus.

\textbf{20.} La signature G en T a prouvé que les \omterm{def:g11:mutations:mutation}{mutations} de la
\omterm{def:g11:cancer:cancer}{tumeur} venaient de la fumée de tabac~; le \omterm{def:g11:cancer:cancer}{cancer} a été fait par un
oncogène activé et les deux \omterm{def:g10:universal-dna:gene}{allèles} de p53 perdus~; ne pas fumer
l'aurait évité.
\end{solution}
